\section{Question and contribution}
If reversing q excitation leaves the effective machine unchanged, direct-axis
flux should retain its value and quadrature-axis flux should reverse sign.
This provides a reference relationship without knowing either inductance.
The research question is which flux components this relationship detects,
and what correction can be justified from reflection alone.

Construction symmetries in energy-based motor models explain the physical
origin of reflection \cite{jebai2014}; nonlinear flux models distinguish
parity from reciprocal differential coupling \cite{sun2020}. Here we isolate
the projection implied by parity, rather than derive another magnetic model.
The contribution is an explicit closest-pair operator, an evidence boundary
for its removed component, and a measured-coordinate demonstration on a
portable PSM fixture. General causal error separation belongs to the companion
identifiability study \cite{companion-diagnostics}.
The algebra applies wherever the declared reflection and matched state hold;
experimental interpretation additionally depends on sampling and coordinates.
